\documentclass[11pt]{amsart}
\usepackage[margin=30mm]{geometry}
\usepackage{amsmath,amssymb,amsthm,mathtools,microtype,hyperref}
\hypersetup{colorlinks=true,linkcolor=blue,citecolor=blue,urlcolor=blue}
\title{Supplement: The Even-Degree Case of Carlet's Conjecture}
\author{Wang Quanzhi}
\date{November 2026}
\newtheorem{theorem}{Theorem}
\newtheorem{lemma}{Lemma}
\begin{document}
\maketitle

\noindent\textbf{Relation to the accompanying proof candidate.} This supplement proves the even-degree statement by a finite-field counting and subspace-polynomial argument. It is a separate route from the reciprocal-geometric argument proposed for odd degree in the accompanying manuscript. The odd-degree argument remains a proof candidate requiring independent specialist review; the existence of this even-degree proof does not certify that argument.

\section{Statement and coefficient criterion}
Let $K$ be a finite field of characteristic $2$ and order $q=2^n$. For an $\mathbb F_2$-subspace $E\subseteq K$, put
\[
L_E(T)=\prod_{x\in E}(T+x)=T^{2^d}+c_{d-1}T^{2^{d-1}}+\cdots+c_1T^2+c_0T,
\qquad d=\dim_{\mathbb F_2}E.
\]
The reciprocal-sum criterion says
\[
\sum_{x\in E\setminus\{0\}}x^{-1}=0\quad\Longleftrightarrow\quad c_1=0.
\]
Indeed, the coefficient of $T^2$ in $L_E(T)$ is the product of the nonzero elements of $E$ times the reciprocal sum; that product is nonzero.

\begin{theorem}
If $n\ge4$ is even, then for every integer $k$ with $2\le k\le n-2$, there is a $k$-dimensional $\mathbb F_2$-subspace $E\subseteq K$ such that $\sum_{x\in E\setminus\{0\}}x^{-1}=0$.
\end{theorem}

\section{Base dimensions two and three}
For even $n$, the field $K$ contains its unique subfield $\mathbb F_4=\{0,1,\omega,\omega^2\}$, where $\omega^2+\omega+1=0$. Then
\[
\sum_{x\in\mathbb F_4^\times}x^{-1}=1+\omega+\omega^2=0,
\]
so dimension two is settled.

Now suppose that $n\ge6$ is even. Choose uniformly among the three-dimensional subspaces containing $1$, and let $p_3(0)$ be the proportion for which the $T^2$ coefficient of $L_E$ is zero. Write $\operatorname{Tr}:K\to\mathbb F_2$ for the field trace and
\[
\mathrm{Kl}_n=\sum_{v\in K^\times}(-1)^{\operatorname{Tr}(v+v^{-1})}.
\]
For later use, let $p_2(v)$ denote the proportion of two-dimensional subspaces containing $1$ whose $T^2$ coefficient is $v$. The exact one-step subspace-extension count, with the even-degree trace term retained, is
\begin{align}
(q-4)p_3(0)&=1-\frac{2}{q-2}+3\sum_{v\in K^\times}(-1)^{\operatorname{Tr}(v^{-1})}p_2(v),\label{eq:rec}\\
(q-2)p_2(v)&=1+(-1)^{\operatorname{Tr}(v)}-2\mathbf 1_{v=1} \qquad(v\ne0).\label{eq:p2}
\end{align}
Here $\mathbf 1_{v=1}$ is $1$ when $v=1$ and $0$ otherwise. For completeness, let $N_d(c)$ be the number of all $d$-subspaces whose $T^2$ coefficient is $c$, let $H_d(c)$ count those that also contain $1$, and let $A_d$ be the total number of $d$-subspaces containing $1$. The raw Artin--Schreier fiber count, obtained by applying the character identity for the fibers of $x\mapsto x^2+x$ to each $d$-subspace, is
\[
2^d H_{d+1}(b)=N_d(b+1)+\sum_{w\ne0}(-1)^{\operatorname{Tr}(w)}H_d((b+1)/w),\qquad b\ne1.
\]
Incidence counting gives $(2^d-1)N_d(c)+H_d(0)=A_d$ for $c\ne0$, and counting adjacent dimensions gives $A_{d+1}/A_d=(q-2^d)/(2^d(2^d-1))$. Dividing the raw identities by the appropriate totals yields \eqref{eq:rec} at $d=2,b=0$. In \eqref{eq:p2}, the two-dimensional coefficient count is explicit; the correction at $v=1$ removes the repeated line, and $\operatorname{Tr}(1)=0$ because $n$ is even.

Multiply \eqref{eq:p2} by $(-1)^{\operatorname{Tr}(v^{-1})}$ and sum over $v\ne0$. The nontrivial additive-character sum over $K^\times$ is $-1$, the product of the two characters is $(-1)^{\operatorname{Tr}(v+v^{-1})}$, and the exceptional $v=1$ term contributes $-2$. Consequently
\[
(q-2)\sum_{v\in K^\times}(-1)^{\operatorname{Tr}(v^{-1})}p_2(v)=\mathrm{Kl}_n-3.
\]
Substituting this into \eqref{eq:rec} and simplifying gives the exact identity
\[
(q-2)(q-4)p_3(0)=q-13+3\mathrm{Kl}_n.
\]

Consider the elliptic curve $C:y^2+xy=x^3+1$. Its discriminant is $1$, so it is nonsingular in characteristic $2$. For each $x\ne0$, the equation in $y$ has $1+(-1)^{\operatorname{Tr}(x+x^{-1})}$ solutions; at $x=0$ it has one solution, and the projective point at infinity contributes one more. Hence
\[
\#C(K)=q+1+\mathrm{Kl}_n.
\]
The Hasse bound for elliptic curves (see \cite{Silverman}) yields $|\mathrm{Kl}_n|\le2\sqrt q$. Since $q\ge64$,
\[
q-13+3\mathrm{Kl}_n\ge q-13-6\sqrt q=s^2-6s-13\ge3>0,
\qquad s=\sqrt q\ge8.
\]
Also $(q-2)(q-4)>0$. Thus $p_3(0)>0$, which proves existence of a three-dimensional zero-sum subspace for every even $n\ge6$. When $n=4$, the allowed range contains only $k=2$, already handled above.

\section{Dimension extension}
\begin{lemma}
If zero-sum subspaces of dimensions $a$ and $b$ exist in $K$ and $ab<n$, then a zero-sum subspace of dimension $a+b$ exists.
\end{lemma}
\begin{proof}
Let $E,F$ be such subspaces and let $I=L_E(K)$. The map $L_E:K\to K$ is $\mathbb F_2$-linear with kernel $E$, so $K/I$ has dimension $a$. Consider the linear map
\[
K\longrightarrow\operatorname{Hom}_{\mathbb F_2}(F,K/I),\qquad
\lambda\longmapsto(x\mapsto\lambda x\bmod I).
\]
Its target has dimension $ab<n$, so its kernel contains a nonzero $\lambda$. Hence $\lambda F\subseteq I$. Put $G=L_E^{-1}(\lambda F)$. The kernel-image dimension formula gives $\dim G=a+b$, and the root sets give $L_G=L_{\lambda F}\circ L_E$. Scaling preserves the vanishing of the $T^2$ coefficient. For linearized polynomials $A,B$, the $T^2$ coefficient of $A\circ B$ is
\[
A'(0)\,[T^2]B+[T^2]A\,( [T]B)^2.
\]
Both terms vanish here because $E$ and $\lambda F$ are zero-sum. Therefore $G$ is zero-sum.
\end{proof}

\section{Complementary dimensions and conclusion}
For any $E$, set $I=L_E(K)$. Since $L_E$ has kernel $E$, $\dim I=n-\dim E$. Moreover,
\[
L_I\circ L_E=T^q-T,
\]
because both sides are monic degree-$q$ polynomials with every element of $K$ as a root. Comparing the $T^2$ coefficients shows that, for $q>2$, $E$ is zero-sum if and only if $I$ is zero-sum. Thus existence in dimension $k$ is equivalent to existence in dimension $n-k$.

We prove existence first for $2\le d\le n/2$. Dimensions two and three are the base cases. For $d\ge4$, apply the dimension-extension lemma to dimensions $d-2$ and $2$. Its hypothesis is $2(d-2)<n$, which follows from $d\le n/2$. Induction proves the claim through the lower half. For any $2\le k\le n-2$, use this result directly if $k\le n/2$; otherwise $2\le n-k<n/2$, and the complementary-dimension equivalence supplies dimension $k$. This proves the theorem for every finite field of order $2^n$, without restricting to a particular field model.

\begin{thebibliography}{1}
\bibitem{Silverman} J. H. Silverman, \emph{The Arithmetic of Elliptic Curves}, second edition, Graduate Texts in Mathematics 106, Springer, 2009, Chapter V, Theorem 1.1.
\end{thebibliography}

\end{document}
